% Source-derived chapter 04: Properties of $t$-birationally liftable Galois sections
% Original source: birational-section-conjecture-strong-birationality
\chapter{Properties of $t$-birationally liftable Galois sections}
\label{birational-section-conjecture-strong-birationality-chapter-04}
\source{birational-section-conjecture-strong-birationality}

\begin{remark}[Source section]
\label{birational-section-conjecture-strong-birationality-chapter-04-source}
\source{birational-section-conjecture-strong-birationality}
\source{birational-section-conjecture-strong-birationality}
This chapter reproduces the corresponding section of the retrieved
LaTeX source, including its definitions, statements, examples, and
proofs. It has not yet been translated into Lean.
\notready
\end{remark}

% The source section title was: Properties of $t$-birationally liftable Galois sections
\label{tprop}
\source{birational-section-conjecture-strong-birationality}

%~ We call a morphism $X\to C$ a \emph{family of curves} if there exists a smooth, projective morphism $\bar{X}\to C$ with connected fibers of dimension $1$ and a divisor $D\s C$ finite étale over $C$ such that $X=\bar{X}\setminus D$. A family of curves is a geometric fibration, thus it has a relative étale fundamental gerbe $\Pi_{X/C}$.

%~ \subsection{Specializations of $t$-birational sections}

In order to work with $t$-b.l. Galois sections, we need to prove a number of facts about them.

\begin{lemma}
\source{birational-section-conjecture-strong-birationality}
\notready
    Let $R$ be a DVR, $X\to \spec R$ a family of curves and $p\in X$ a closed point in the closed fiber. There exists a non-empty divisor $D\subset X$ finite étale over $R$ with $p\in D$.
	\begin{proof}
		%~ Let $K$, $k$ be the fraction and residue fields of $R$. First, let us show that it is enough to prove the statement in the case in which $p$ is $k$-rational. Let $k'$ be the residue field of $p$, since we are in characteristic $0$ we have $k'\simeq k[x]/\bar{q}(x)$ for some irreducible, separable polynomial $\bar{q}\in k[x]$, let $q\in R[x]$ any lifting of $\bar{q}$. By construction, $R'=R[x]/q(x)$ is finite étale over $R$ and has only one closed point, hence it is a DVR. If $D'\subset X_{R'}$ is finite étale over $R'$, it is finite étale over $R$, too, hence the image of $D'$ in $X$ is finite étale over $R$. 
		
		%~ Assume that $p$ is $k$-rational. By definition, $X\to\spec R$ has the form $\bar{X}\setminus D_{0}\to\spec R$ with $\bar{X}\to\spec R$ smooth, projective, with geometrically connected fibers of dimension $1$, and $D_{0}\subset\bar{X}$ a divisor finite étale over $R$. Choose a projective embedding $\bar{X}\subset\P^{n}_{R}$.
		
		%~ A generic hyperplane $H_{k}$ of $\P^{n}_{k}$ passing through $p$ has normal crossings with $\bar{X}_{k}$ and does not intersect $D_{0,k}$. Let $q\in R[t_{0},\dots,t_{n}]$ be a linear form whose image in $k[t_{0},\dots,t_{n}]$ is an equation for $H_{k}$, then $q$ defines a closed subscheme $H\subset\P^{n}_{R}$ which is a relative hyperplane. The fact that $H_{k}\cap D_{0,k}=\emptyset$ implies that $H\cap D_{0}=\emptyset$ since everything is proper over $R$, and the fact that $H_{k}$ has normal crossings with $\bar{X}_{k}$ implies that $D=H\cap X=H\cap\bar{X}$ is finite étale over $R$.
		
		
		
		%~ \newpage
		
		Let $K,k$ be the fraction and residue fields of $R$. If $p$ is $k$-rational, this essentially follows from Bertini's theorem by taking a generic hyperplane section passing through $p$. If $p$ is not $k$-rational, though, this is more subtle. We are going to use a Bertini-like argument to reduce to the case of $\P^{1}_{R}$, where we are able to make an explicit construction.
		
		By definition, $X\to\spec R$ has the form $\bar{X}\setminus D_{0}\to\spec R$ with $\bar{X}\to\spec R$ smooth, projective, with geometrically connected fibers of dimension $1$, and $D_{0}\subset\bar{X}$ a divisor finite étale over $R$. Choose a projective embedding $\bar{X}\subset\P^{n}_{R}$ and let $p_{1},\dots,p_{r}$ be the geometric points over $p$.
			
		A generic $(n-2)$-dimensional linear subspace $L_{k}$ of $\P^{n}_{k}$ has the following properties
		\begin{itemize}
			\item $L_{k}\cap \bar{X}_{k}=\emptyset$,
			\item for every $i$ the hyperplane $H_{i}\subset\P^{n}_{\bar{k}}$ spanned by $L_{k}$ and $p_{i}$ has transversal intersections with $\bar{X}_{\bar{k}}$,
			\item $\forall i:H_{i}\cap D_{0}=\emptyset$.
		\end{itemize}
		
		Let $\bar{q}_{0},\bar{q}_{1}\in k[t_{0},\dots,t_{n}]$ two $k$-linear forms which are equations for $L_{k}$, choose $q_{0},q_{1}\in R[t_{0},\dots,t_{n}]$ $R$-linear forms which lift $\bar{q}_{0},\bar{q}_{1}$ and let $L\subset\P^{n}_{R}$ be their vanishing locus, the fact that $L_{k}\cap \bar{X}_{k}=\emptyset$ implies that $L\cap \bar{X}=\emptyset$ since everything is proper over $R$.
		
		The graded homomorphism $R[x_{0},x_{1}]\to R[t_{0},\dots,t_{n}]$, $x_{i}\mapsto q_{i}$ defines a morphism 
		\[\P^{n}_{R}\setminus L=\proj R[t_{0},\dots,t_{n}]\setminus L\to\P^{1}_{R}=\proj R[x_{0},x_{1}].\]
		We have simply constructed the projection centered in $L$: we can do this even though we don't have a base field.
		
		Since $X\cap L=\emptyset$, we get a morphism $f:\bar{X}\to\P^{1}_{R}$ over $R$ with $\bar{X}_{k}\to\P^{1}_{k}$ surjective. The morphism $f$ is flat since it is finite and dominant, and both $\bar{X}$ and $\P^{1}_{R}$ are regular. Moreover, $f$ is unramified, and hence étale, at $\spec k(X_{k})\subset X$: notice that the diagram
		\[\begin{tikzcd}
			\bar{X}_{k}\rar\dar	&	\P^{1}_{k}\rar\dar	&	\spec k\dar	\\
			\bar{X}\rar			&	\P^{1}_{R}\rar		&	\spec R
		\end{tikzcd}\]
		is cartesian, hence $\spec k(\P^{1}_{k})\times_{\P^{1}_{R}}\bar{X}=\spec k(\bar{X}_{k})$ is reduced. By purity, it follows that the branch locus of $f$ is the closure of a finite number of closed points of $\P^{1}_{K}$. Let $K''/K$ be some finite extension and $H_{K''}\subset \P^{n}_{K''}$ a hyperplane containing $L_{K''}$ and corresponding to a branch point in $\P^{1}_{K}$; equivalently, $H_{K''}$ does not have transversal intersections with $X_{K''}$. If $H_{k''}\subset \P^{n}_{k''}$ is a specialization of $H_{K''}$ in some finite extension $k''$ of $k$, then $H_{k''}$ does not have transversal intersections with $X_{k''}$, it follows that $H_{k''}$ does not contain $p_{i}$ for every $i$. This implies that $H_{k''}$ corresponds to a point of $\P^{1}(k'')$ different from $f(p)$, hence $f(p)$ is not a branch point for $f$.
		
		Let $U_{0}\subset\P^{1}_{R}$ be the open subset where $f$ is étale, i.e. the complement of the image of the ramification locus in $\bar{X}$, and $U=U_{0}\setminus f(D_{0})$. Since $f(p)$ is not a branch point and $H_{i}\cap D_{0}=\emptyset$ for every $i$ we get that $f(p)\in U$. Furthermore, $f^{-1}(U)\to U$ is finite étale by construction. It is then enough to find a divisor of $U$ containing $f(p)$ and which is finite étale over $R$. 
		 %~ fact that $H_{i}$ has normal crossing with $\bar{X}_{\bar{k}}$ for every $i$ implies that $f(p)\in U_{0}$, while 
		
		%~ the fact that $H_{i}\cap D_{0}=\emptyset$ implies that $f(p)\in U$. Furthermore, $f^{-1}(U)\to U$ is finite étale by construction. It is then enough to find a divisor of $U$ containing $f(p)$ and which is finite étale over $R$.
		
		Since we are in characteristic $0$, then $k(p)\simeq k[x]/\bar{q}(x)$ for some irreducible, separable polynomial $\bar{q}\in k[x]$, let $q\in R[x]$ be any lifting of $\bar{q}$. By construction, $R'=R[x]/q(x)$ is finite étale over $R$ and has only one closed point, hence it is a DVR. Let $K'$ be the fraction field of $R'$, there are infinitely many points of $\P^{1}(K')$ which specialize to $p\in\P^{1}(k(p))$, hence we may choose one in $U(K')$; denote by $E\subset \P^{1}_{R}$ its closure, we have an extension $\spec R'\to E\subset\P^{1}_{R}$. By construction, $E$ is finite and flat over $R$ of degree equal to $[K':K]$. Furthermore, $E$ contains $p$, which satisfies $[k(p):k]=[K':K]$. Because of this, $E$ contains only $p$ and a generic point contained in $U$, hence $E\subset U$ and $\spec R'\to E$ is surjective. Since $R'$ is étale over $R$ we get that $E$ is étale over $R$ as well (we have actually proved that $\spec R'\to E$ is an isomorphism, though we don't need this).
  
%  The extension $\spec R'\to U$ is a closed embedding by degree reasons and defines the desired divisor.
		
		Notice that if we try to work with $R'$-sections of $\bar{X}$ we run into problems: it might be that $X(K')$ does not contain points which specialize to $p$, and if we try to enlarge $R'$ again we might get a semi-local ring with several closed points which we are not able to control simultaneously. Making a Bertini-like argument over $\bar{X}_{R'}$ runs into problems too: while we may find the desired divisor étale over $R'$, there is no guarantee that its image in $\bar{X}$ is étale over $R$.
		
		
		
		
		
		 %~ The fact that the hyperplanes $H_{i}$ have normal crossing with $\bar{X}_{k'}$ imply that $f(p)\in U$, while the fact that $H_{i}\cap D_{0}=\emptyset$ implies that $f(p)\not\in f(D_{0})$.\todo{BOOKMARK}
			
			%~ \newpage


            %~ By hypothesis, $R$ is a localization of a regular $k_{0}$-algebra of finite type. Let $k$ be the residue field of $R$, it is a finite extension of $k_{0}$. Let $k'/k_{0}$ be a Galois extension such that $k(p)/k_{0}$ splits completely over $k'$. Fix a preferred embedding $k\subset k'$, let $p_{1},\dots,p_{r}\in X_{k}(k')$ be the $r=[k(p):k]$ $k'$-rational points over $p$, $\gal(k'/k)\subset\gal(k'/k_{0})$ acts transitively on them.

			%~ Define $R'=R\otimes_{k_{0}}k'$, it is a semi-local ring finite étale over $R$ whose residue fields are $[k:k_{0}]$ copies of $k'$ and there is an induced action of $\gal(k'/k_{0})$ on $R'$.
			
			
			 %~ The preferred embedding $k\subset k'$ defines a preferred closed point of $\spec R'$, 
			
			%~ Each point $p_{i}\in X(k')$ induces a point of $X_{R'}(k')$ which we still call $p_{i}$, $\gal(k'/k)\subset\gal(k'/k_{0})$ acts transitively on them. It is sufficient to find a $\gal(k'/k)$ invariant divisor of $X_{R'}$ which contains the points $p_{i}$ and is finite étale over $R'$.
            

            %~ By definition, $X\to\spec R$ has the form $\bar{X}\setminus D_{0}\to\spec R$ with $\bar{X}\to\spec R$ smooth, projective, with geometrically connected fibers of dimension $1$, and $D_{0}\subset\bar{X}$ a divisor finite étale over $R$. Choose a projective embedding $\bar{X}\subset\P^{n}_{R}$.
		

		
		%~ Let $k_{R}$ be the residue field of $R$, we have embeddings $k\subset k_{R}\subset k(p)\subset k'$. A generic $(n-2)$-dimensional linear subspace $L_{k}$ of $\P^{n}_{k}$ has the following properties
		%~ \begin{itemize}
			%~ \item $L_{k}\cap \bar{X}_{k}=\emptyset$,
			%~ \item for every $i$ the hyperplane $H_{i}\subset\P^{n}_{k'}$ spanned by $L_{k}$ and $p_{i}$ has normal crossing with $\bar{X}_{k'}$,
			%~ \item $\forall i:H_{i}\cap D_{0}=\emptyset$,
			%~ \item $\forall i\neq j: H_{i}\neq H_{j}$, i.e. $p_{i}\notin H_{j}$ and $H_{i}\cap H_{j}=L_{k}$.
		%~ \end{itemize} 
		
		%~ Choose such an $L_{k}$, let $L_{K}$ be any $n-2$-dimensional linear subspace of $\P^{n}_{K}$ which specializes to $L_{k}$. By construction, $L_{K}\cap \bar{X}_{K}=\emptyset$. Choose $q_{1}\in X_{K}(\bar{K})$ any point which specializes to $p_{1}$, there is a finite number of Galois conjugates of $q_{1}$ and each \todo{ARGHHHH}
		
		
		 %~ Let $R'$ be a DVR over $R$ with fraction and residue fields $K',k'$ such that $k(p)$ splits completely in $k'$, we may thus regard the points $p_{1},\dots,p_{r}$ as $k'$-rational point. Let $H_{i}\subset\P^{n}_{k'}$ be the hyperplane spanned by $L_{k}$ and $p_{i}$, by construction $H_{i}\neq H_{j}$ for $i\neq j$. 
		
		
		
		
		%~ We claim that the fibers of the specialization map $\P^{n}(K')\to\P^{n}(k')$ are Zariski-dense. It is enough to show this for the point $(1:0:0:\dots:0)$, whose inverse images are the points $(1:a_{1}:\dots:a_{n})$ such that $a_{i}$ is in the maximal ideal of $R'$ for every $i\neq 0$. If $h\in K[t_{0},\dots,t_{n}]$ is a non-zero, homogeneous polynomial which vanishes on all such points, up to permuting the indexes we get that $h(1,t_{1},0,\dots,0)$ is a non-zero polynomial in one variable with infinite roots, which is absurd.
		
		%~ The hyperplanes of $\P^{n}_{\bar{K}}$ containing $L_{K}$ are parametrized by $\P^{1}(\bar{K})$. Let us say that an hyperplane of $\P^{n}_{\bar{K}}$ containing $L_{K}$ is bad if it intersects $D_{0}$ or if it has not normal crossing with $X_{K}$
		
		%~ Let $Z\subset\P^{n}_{\bar{K}}$ be the union of the $\gal(\bar{K}/K)$-conjugates of the hyperplanes containing $L_{k}$
		
		%~ Let $U\subset G_{R}(k)$ be a Zairski-open subset of linear subspaces with such properties, we may assume $U\neq\emptyset$ since $k$ is infinite.
		
		%~ we have $L_{k}\cap \bar{X}_{k}=\emptyset$ and, for every $i=1,\dots,r$, the hyperplane $H_{i}\subset\P^{n}_{\bar{k}}$ spanned by $L$ and $p_{i}$ has normal crossing with $X_{\bar{k}}$ and does not intersect $D_{0}$.
		
		%~ For a generic $(n-2)$-dimensional linear subspace $L_{K}$ of $\P^{n}_{K}$, we have $L_{K}\cap \bar{X}_{K}=0$. Let $V\subset G_{R}(K)$ be a Zariski open subset of such linear subspaces.
		
		%~ Consider the specialization map $s:G_{R}(K)\to G_{R}(k)$: it is \emph{not} Zariski-continuous, but it does map proper, closed subsets to non-dense subsets. To check this, consider some projective embedding $G_{R}\subset\P^{N}_{R}$, a proper closed subscheme $Z\subset G_{K}$ and a homogeneous polynomial $p\in K[t_{0},\dots,t_{N}]$ which vanishes on $Z$ but not on $G_{K}$. Up to multipling $p$ by a suitable scalar, we may assume that $p\in R[t_{0},\dots,t_{N}]$ and that the reduction $\bar{p}$ of $p$ in $k[t_{0},\dots,t_{N}]$ is non-zero. Then $s(Z(K))\subset \{\bar{p}=0\}\cap G_{k}(k)$, which is a proper closed subset of $G_{k}(k)$. In particular, we get that inverse images of dense subsets are dense.
		
		%~ Since $U$ is dense in $G_{R}(k)$, it follows that the inverse image $s^{-1}(U)$ of $U$ in $G_{R}(K)$ is dense and we may choose $L_{K}\in V\cap s^{-1}(U)$. Let $L_{R}\subset\P^{n}_{R}$ be the closure of $L_{K}$, the closed fiber $L_{k}$ of $L_{R}$ is an element of $U\subset G_{R}(k)$. Let $R'$ be a DVR over $R$ with fraction and residue fields $K',k'$ such that $k(p)$ splits completely in $k'$, we may thus regard the points $p_{1},\dots,p_{r}$ as $k'$-rational point. Let $H_{1},\dots,H_{l}\subset\P^{n}_{k'}$ be the hyperplanes spanned by $L_{k'}$ plus one of the points $p_{1},\dots,p_{r}$: notice that 
		
		%~ \newpage
		
		
		 %~ Let $k'$ be the residue field of $p$, since we are in characteristic $0$ is a separable extension of $k$. There exists a DVR $R'$ unramified over $R$ such that \todo{BOOKMARK WAAAAA}
		
		
		%~ By definition, $X\to\spec R$ has the form $\bar{X}\setminus D_{0}\to\spec R$ with $\bar{X}\to\spec R$ smooth, projective, with geometrically connected fibers of dimension $1$, and $D_{0}\subset\bar{X}$ a divisor finite étale over $R$. Choose a projective embedding $\bar{X}\subset\P^{n}_{R}$ and let $K$ be the fraction field of $R$, we may assume that $\bar{X}_{K}\subset\P^{n}_{K}$ is not contained in any hyperplane. A generic hyperplane of $\P^{n}_{K}$ has normal crossings with $X_{K}$ and does not intersect $D_{0}$, the same is true over the closed fiber.
		
		%~ Consider the dual projective bundle $P\simeq\P^{n}_{R}\to\spec R$ parametrizing hyperplanes and let $k$ be the residue field, there is a specialization map $s:P(K)\to P(k)$. Let $U\subset P(k)$ be a Zariski open subset of hyperplanes which have normal crossing with $X_{k}$ and avoid $D_{0}$, we may assume that $U$ is non-empty since $k$ is infinite (we are in characteristic $0$). The specialization map $s$ is \emph{not} Zariski-continuous, but it maps proper closed subsets to proper closed subsets: this follows from the fact that we can multiply every non-zero homogeneous polynomial with coefficients in $K$ by a scalar so that the new coefficients are in $R$ and the reduction to $k$ is non-zero. Because of this, the inverse image $s^{-1}(U)\subset P(K)$ is dense. Choose an hyperplane $H_{K}\in s^{-1}(U)\subset P(K)$ which has normal crossing with $X_{K}$ and avoids $D_{0}$, since $H_{K}\in s^{-1}(U)$ the same is true for the specialization $H_{k}$ of $H_{K}$, i.e. the closed fiber of the closure $H_{R}\subset\P^{n}_{R}$ of $H_{K}$. 
		
		%~ Choose $D$ as the intersection $H_{R}\cap\bar{X}$ in $\P^{n}_{R}$. By construction, $D\cap D_{0}=\emptyset$, $D\to\spec R$ is finite and $D_{k}=H_{k}\cap X_{k}$ is reduced, hence $D\to\spec R$ is finite étale.
		
		
		
		
		
		%~ Let $\bar{X}\to \spec R$ be a family of projective curves with $X\s\bar{X}$. Choose any generically finite rational map $f:\bar{X}\dashrightarrow\P^{1}_{R}$ over $R$ such that the special fiber is not a ramification divisor and let $f':\bar{X}'\to\P^{1}_{R}$ be a resolution. Let $k$ be the residue field, there exists an open subset $U\s X\s\bar{X}$ with $U_{k}\neq\emptyset$ and such that $U\to\P^{1}_{R}$ is an étale morphism. Let $Z=\bar{X}'\setminus U$ be the complement and write $V=\bar{X}'\setminus f^{'-1}(f'(Z))$, we have that $V$ is contained in $U$ and $V=f^{'-1}(f'(V))\to f'(V)$ is finite étale. Since $k$ has characteristic $0$ we may chose a rational point $p$ in the special fiber $f'(V)_{k}\s\P^{1}_{k}$. Since the fibers of the homomorphism $R^{*}\to k^{*}$ are infinite, we may extend $p$ to a section $s:\spec R\to f'(V)\s\P^{1}_{R}$. The inverse image $f^{-1}(s(\spec R))\s X$ is finite étale over $R$.
	\end{proof}
\end{lemma}

\begin{corollary}
\source{birational-section-conjecture-strong-birationality}
\notready\label{cut}
\source{birational-section-conjecture-strong-birationality}
	Let $R$ be a DVR, $X\to \spec R$ a family of curves. There exists a direct system $(D_{i})_{i}$ of divisors finite étale over $R$ such that $X_{k}\setminus\bigcup_{i}D_{i}$ contains only the generic point of $X_{k}$.
\end{corollary}


\begin{lemma}
\source{birational-section-conjecture-strong-birationality}
\notready\label{spebir}
\source{birational-section-conjecture-strong-birationality}
	Specializations of b.l. sections are b.l.
	\begin{proof}
		Let $X\to \spec R$ be a family of curves with $R,K,k$ as above and $c\in\spec R$ the closed point. If $z\in\Pi_{X_{K}/K}(K)$ is a generic, b.l. Galois section, we have a specializing loop $\gamma_{z}:S^{1}_{c}\to\Pi_{X_{k}/k}$. 
		%~ We want to prove that the sections of $\Pi_{X_{p}/k(p)}(k(p))$ in the essential image of $\gamma_{z}(k(p))$ are birational.
		
		Let $D_{i}\s X$ be a direct system of divisors as in \autoref{cut}. We have that $X_{i}=X\setminus D_{i}$ is a family of curves for every $i$, write $X_{\infty}=\projlim_{i}X_{i}$ and $\Pi_{X_{\infty}/R}=\projlim_{i}\Pi_{X_{i}/R}$. The closed fiber of $X_{\infty}$ is naturally isomorphic to $\spec k(X_{k})$.
		
		Since $z$ is b.l. and $\spec K(X)\s X_{\infty,K}$, there exists a lifting $z'\in\Pi_{X_{\infty,K}/K}(K)$ of $z$ to $X_{\infty,K}$. By a limit argument, the specializing loops of $z'$ in $\Pi_{X_{i}/R}$ induce a specializing loop $\gamma_{z'}:S^{1}_{c}\to \Pi_{X_{\infty,k}/k}=\Pi_{k(X_{k})/k}$, and the composition of $\gamma_{z'}$ with $\Pi_{k(X_{k})/k}\to \Pi_{X_{k}/k}$ is isomorphic to $\gamma_{z}$. Hence, every specialization of $z$ lifts to $\spec k(X_{k})$.
	\end{proof}
\end{lemma}

\begin{corollary}
\source{birational-section-conjecture-strong-birationality}
\notready\label{tbb}
\source{birational-section-conjecture-strong-birationality}
	 $t$-b.l. sections are b.l. 
	 \begin{proof}
		Let $X$ be a curve over a field $k$ and $s\in\Pi_{X/k}(k)$ a $t$-b.l. Galois section. Write $R=k[t]_{(t)}$, we have that $s_{k(t)}$ is a generic section of $\Pi_{X_{R}/R}$ and $s$ is a specialization of $s_{k(t)}$. By hypothesis $s_{k(t)}$ is b.l., hence its specialization $s$ is b.l. by \autoref{spebir}. 
	 \end{proof}
\end{corollary}

\begin{lemma}
\source{birational-section-conjecture-strong-birationality}
\notready\label{spetbir}
\source{birational-section-conjecture-strong-birationality}
	Specializations of $t$-b.l. sections are $t$-b.l.
	\begin{proof}
		Let $X\to \spec R$ be as above, $z\in\Pi_{X_{K}/K}(K)$ a generic, $t$-b.l. Galois section, $s\in\Pi_{X_{k}/k}(k)$ a specialization. Let $R'$ be the local ring of the generic point of the divisor $\A^{1}_{k}\s\A^{1}_{R}$. We have that $R'$ is a DVR with fraction field $K'=K(t)$ and residue field equal $k'=k(t)$. Consider the family of curves $X_{R'}\to \spec R'$, then $z_{K(t)}=z_{K'}\in \Pi_{X_{K'}/K'}(K')$ is by hypothesis a b.l. Galois section. We have that $s_{k(t)}=s_{k'}$ is a specialization of $z_{K'}$ and thus it is b.l. by \autoref{spebir}.
	\end{proof}
\end{lemma}

%~ \subsection{Lifting $t$-birational sections}

\begin{lemma}
\source{birational-section-conjecture-strong-birationality}
\notready\label{morb}
\source{birational-section-conjecture-strong-birationality}
	Let $f:Y\to X$ be a dominant morphism of curves over a field $k$ and $s\in\Pi_{Y/k}(k)$ a Galois section. If $s$ is b.l. then $f(s)\in\Pi_{X/k}(k)$ is b.l. If $f$ is finite étale, the converse holds.
	\begin{proof}
		The first statement is obvious. If $f$ is finite étale, the second statement follows from the fact that $\gal(k(Y))=\pi_{1}(Y)\times_{\pi_{1}(X)}\gal(k(X))$.
	\end{proof}
\end{lemma}

\begin{corollary}
\source{birational-section-conjecture-strong-birationality}
\notready\label{mortb}
\source{birational-section-conjecture-strong-birationality}
	Let $f:Y\to X$ be a dominant morphism of curves over a field $k$ and $s\in\Pi_{Y/k}(k)$ a Galois section. If $s$ is $t$-b.l. then $f(s)\in\Pi_{X/k}(k)$ is $t$-b.l. If $f$ is finite étale, the converse holds.\qed
\end{corollary}

We say that a curve $X$ is \emph{parabolic} if it has genus $0$ and the degree of $\bar{X}\setminus X$ is at most $2$. If $X$ is proper of genus $0$, then there is a unique Galois section, and it is geometric if and only if $X=\P^{1}$. If $X$ is parabolic affine, every Galois section of $X$ is cuspidal (the only non-trivial case is settled by \cite[Theorem A]{sch15}). Because of this, we can always assume that $X$ is non-parabolic.

If $X$ is a curve over a number field $k$ and $s$ is a b.l. section, by Koenigsmann's results \cite{koe05}, \cite[Proposition 1]{sti15} for every finite place $\nu$ of $k$ there exists a local point $x_{\nu}\in\bar{X}(k_{\nu})$ such that $s_{k_{\nu}}$ is associated with $x_{\nu}$, i.e. if $x_{\nu}\in X$ then $s_{\nu}$ is the geometric section associated with $x_{\nu}$, otherwise $s_{\nu}$ is one of the cuspidal sections of the packet associated with $x_{\nu}\in \bar{X}\setminus X$.

If $X$ is non-parabolic, the local point $x_{\nu}$ associated with $s$ is unique. To check this, we can pass to an étale neighbourhood $Y$ of genus $\ge 1$ and use the injectivity of the section map \cite[Proposition 75]{sti13} on $\bar{Y}_{k_{\nu}}$.

%~ \begin{lemma}
	%~ Let $k$ be a finite extension of $\Q_{p}$, $X$ a curve over $k$ with good reduction. Suppose we have a morphism $h:B\hz(1)\to\Pi_{X/k}$ such that there are only finitely many isomorphism classes in the essential image of $B\hz(1)(k)\to\Pi_{X/k}(k)$. Then $h$ is constant.
	%~ \begin{proof}
			%~ If $h$ is not constant, up passing to an étale cover of $X$ we may assume that the composition $h^{\rm ab}:B\hz(1)\to\Pi_{X/k}\to\Pi_{X/k}^{\rm ab}$ is not constant, where $\Pi_{X/k}^{\rm ab}$ is the abelianization.
	%~ \end{proof}
%~ \end{lemma}

\begin{lemma}
\source{birational-section-conjecture-strong-birationality}
\notready\label{nflift}
\source{birational-section-conjecture-strong-birationality}
	Let $X$ be a curve over a number field $k$, $U\s X$ a non-empty open subset and $s\in\Pi_{X/k}(k)$ a $t$-b.l. Galois section. There exists a $t$-b.l. section of $U$ which lifts $s$.
	\begin{proof}
		We may assume that $X$ is non-parabolic and that $s$ is neither geometric nor cuspidal, otherwise this is trivial. 
		
		By hypothesis, $s_{k(t)}$ lifts to a b.l. Galois section $r_{0}$ of $U_{k(t)}$, we are going to show that $r_{0}$ is the base change to $k(t)$ of some Galois section of $U$. Identify $k(t)$ with the fraction field of $\A^{1}$ and choose $c\in \A^{1}$ a closed point with residue field $k'$.
  
        For every finite place $\nu$ of $k$ there is a unique local point $x_{\nu}\in\bar{X}(k_{\nu})$ associated with $s$. Since $s$ lifts to a b.l. section of $U$ by \autoref{tbb} and it is not geometric nor cuspidal, the set of places $\nu$ such that $x_{\nu}\in U$ has Dirichlet density $1$ by \cite[Theorem B]{sti15}, hence we may choose $\nu$ with $x_{\nu}\in U$ and $k'\subset k_{\nu}$. If $r_{0,c}$ is a specialization of $r_{0}$ at $c$, by construction it is b.l. and lifts $s$, hence $r_{0,c,k_{\nu}}$ is geometric associated with $x_{\nu}$. By \autoref{uniqueloop1}, the specializing loop $S^{1}_{c}\to\Pi_{U_{k'}/k'}$ of $r_{0}$ at $c$ is constant and hence $r_{0}$ extends to a section $\spec\O_{c}\to\Pi_{U\times\A^{1}/\A^{1}}$ by \autoref{prorootval}.
		
		Since we can do this for every closed point $c\in\A^{1}$ and we may write $\Pi_{U\times\A^{1}/\A^{1}}$ as a product $\Pi_{U/k}\times_{k}\A^{1}$, by \cite[Corollary A.3]{fed} we get that $r_{0}$ extends to a section $\tilde{r}_{0}:\A^{1}\to\Pi_{U\times\A^{1}/\A^{1}}$. Since $\Pi_{\A^{1}/k}=\spec k$, the composition $\A^{1}\to\Pi_{U\times\A^{1}/\A^{1}}\to\Pi_{U/k}$ factorizes through a section $r:\spec k\to\Pi_{U/k}$ and hence $\tilde{r}_{0}=r_{\A^{1}}$, $r_{0}=r_{k(t)}$. This implies that $r$ is a $t$-b.l. lifting of $s$.
	\end{proof}
\end{lemma}
